\section{Thiết kế và Triển khai Cơ sở dữ liệu}

Đây là nội dung trọng tâm của đồ án. Toàn bộ mục này trình bày từ mô hình dữ liệu logic, quyết định chuẩn hóa, đến schema vật lý PostgreSQL thật kèm mã nguồn SQL, cơ chế toàn vẹn dữ liệu và chiến lược hiệu năng — tương ứng đúng với phạm vi bàn giao "Database làm trọng tâm" mà nhóm đã thống nhất.

\subsection{Danh mục thực thể (Entity Catalog)}
Toàn bộ 17 thực thể được thiết kế từ ERD (Chương 3, mục 3.2.1.5) được triển khai thành bảng vật lý trong PostgreSQL. Mọi thực thể (trừ Tenant) đều có khóa ngoại TenantId để đảm bảo cô lập dữ liệu đa khách thuê.

\begin{longtable}{|p{3.4cm}|p{5.5cm}|p{5cm}|}
    \hline
    \textbf{Thực thể} & \textbf{Mô tả} & \textbf{Thuộc tính chính (rút gọn)} \\
    \hline
    \endfirsthead
    \hline
    \textbf{Thực thể} & \textbf{Mô tả} & \textbf{Thuộc tính chính (rút gọn)} \\
    \hline
    \endhead
    Tenant & Đại diện một cửa hàng/doanh nghiệp thuê hệ thống & TenantId (PK), Name, CreatedAt, Status \\
    \hline
    User & Tài khoản người dùng thuộc một Tenant & UserId (PK), TenantId (FK), Email, PasswordHash, Role \\
    \hline
    Branch & Chi nhánh/kho hàng thuộc một Tenant & BranchId (PK), TenantId (FK), Name, Address \\
    \hline
    Category & Danh mục sản phẩm, phân cấp cha-con & CategoryId (PK), TenantId (FK), ParentCategoryId \\
    \hline
    Brand & Thương hiệu sản phẩm & BrandId (PK), TenantId (FK), Name \\
    \hline
    Product & Sản phẩm gốc, có thể có nhiều biến thể & ProductId (PK), TenantId, CategoryId, BrandId \\
    \hline
    ProductVariant (SKU) & Biến thể cụ thể, đơn vị quản lý tồn kho & SKUId (PK), SKUCode (unique/Tenant), RetailPrice \\
    \hline
    Barcode & Mã vạch gắn với một SKU & BarcodeId (PK), SKUId (FK), Code, Type \\
    \hline
    InventoryTransaction (Ledger) & Sổ cái tồn kho append-only & TransactionId (PK), Type, Quantity, BalanceAfter, CreatedAt \\
    \hline
    StockBalance & Số dư tồn kho hiện tại theo SKU/Branch & SKUId (FK), BranchId (FK), OnHand, Reserved, AvgCost \\
    \hline
    PurchaseReceipt & Phiếu nhập hàng & ReceiptId (PK), BranchId, SupplierInfo, Status \\
    \hline
    PurchaseReceiptItem & Chi tiết dòng hàng trong phiếu nhập & ReceiptItemId (PK), SKUId, Quantity, UnitCost \\
    \hline
    StockTransfer & Phiếu chuyển kho giữa 2 chi nhánh & TransferId (PK), FromBranchId, ToBranchId, Status \\
    \hline
    Stocktake & Phiên kiểm kho & StocktakeId (PK), BranchId, Status \\
    \hline
    StocktakeItem & Chi tiết số đếm thực tế theo SKU & StocktakeItemId (PK), CountedQty, SystemQty, VarianceQty \\
    \hline
    Order & Đơn hàng hợp nhất từ mọi kênh & OrderId (PK), Channel, Status, CreatedAt \\
    \hline
    OrderItem & Chi tiết dòng sản phẩm, snapshot giá vốn & OrderItemId (PK), SKUId, Quantity, CostPrice \\
    \hline
    \caption{Danh mục 17 thực thể triển khai trong schema PostgreSQL}
    \label{tab:entity-catalog}
\end{longtable}

\subsection{Phân tích chuẩn hóa dữ liệu (Normalization Analysis)}

Trước khi viết DDL, nhóm rà soát mô hình dữ liệu theo các dạng chuẩn (Normal Forms) để đảm bảo không dư thừa dữ liệu ngoài chủ đích, đồng thời xác định rõ những chỗ \textbf{chủ động phi chuẩn hóa (denormalization)} vì lý do hiệu năng.

\subsubsection{Dạng chuẩn 1 (1NF)}
Mọi cột đều mang giá trị nguyên tử theo đúng kiểu dữ liệu khai báo (VARCHAR, NUMERIC, TIMESTAMPTZ...). Riêng thuộc tính \texttt{Attributes} của ProductVariant (ví dụ: màu sắc, size) được lưu dưới dạng \texttt{JSONB} — đây không vi phạm 1NF vì PostgreSQL xử lý JSONB như một kiểu dữ liệu nguyên tử ở tầng lưu trữ/index (hỗ trợ index GIN), khác với việc nhét nhiều giá trị phân tách bằng dấu phẩy vào một cột VARCHAR.

\subsubsection{Dạng chuẩn 2 (2NF)}
Các bảng có khóa chính phức hợp đều được kiểm tra để mọi thuộc tính phi khóa phụ thuộc đầy đủ vào toàn bộ khóa chính, không phụ thuộc vào một phần khóa. Ví dụ điển hình: bảng \texttt{StockBalance} có khóa chính phức hợp \texttt{(SKUId, BranchId)} — cả ba thuộc tính \texttt{OnHand}, \texttt{Reserved}, \texttt{AvgCost} đều phụ thuộc vào \emph{cả hai} cột khóa (tồn kho của một SKU là khác nhau tại từng chi nhánh), không có thuộc tính nào chỉ phụ thuộc riêng SKUId hoặc riêng BranchId.

\subsubsection{Dạng chuẩn 3 (3NF)}
Loại bỏ phụ thuộc bắc cầu bằng cách tách các thực thể có vòng đời và tần suất thay đổi khác nhau ra bảng riêng, thay vì lặp lại dữ liệu:
\begin{itemize}
    \item Category và Brand được tách khỏi Product (thay vì lưu tên danh mục/thương hiệu trực tiếp trên từng dòng sản phẩm) để tránh cập nhật hàng loạt khi đổi tên danh mục.
    \item Product và ProductVariant được tách riêng: các thuộc tính chung (tên, danh mục, thương hiệu) nằm ở Product; các thuộc tính riêng theo biến thể (SKU, giá, tồn kho) nằm ở ProductVariant — tránh lặp lại thông tin sản phẩm gốc trên từng biến thể.
\end{itemize}

\subsubsection{Phi chuẩn hóa có chủ đích: StockBalance}
Về mặt lý thuyết, \texttt{StockBalance.OnHand} hoàn toàn có thể được suy ra 100\% bằng cách tính tổng \texttt{InventoryTransaction} theo từng SKU/Branch — nghĩa là StockBalance là dữ liệu \emph{dẫn xuất} (derived data), không phải dữ liệu chuẩn hóa độc lập. Tuy nhiên, nhóm chủ động giữ bảng này như một \textbf{read model} tách biệt (tương tự tư duy CQRS — Command Query Responsibility Segregation), vì hai lý do:
\begin{itemize}
    \item Truy vấn tồn kho khả dụng (\texttt{available = on\_hand - reserved}) là thao tác xảy ra với tần suất rất cao (mỗi lần thêm sản phẩm vào giỏ hàng POS, mỗi lần webhook gửi đơn) — nếu phải \texttt{SUM()} toàn bộ lịch sử Ledger mỗi lần đọc, chi phí sẽ tăng tuyến tính theo số lượng giao dịch lịch sử, không chấp nhận được ở quy mô hàng triệu dòng ledger (NFR-TENANT-02).
    \item StockBalance đóng vai trò là dòng cần \textbf{khóa (lock)} trong cơ chế chống bán vượt tồn (mục~\ref{sec:concurrency-control}) — cần một bảng có đúng 1 dòng/SKU/Branch để \texttt{SELECT...FOR UPDATE}, khác với bảng Ledger vốn có vô số dòng lịch sử.
\end{itemize}
Đánh đổi này được kiểm soát chặt bằng script kiểm tra toàn vẹn dữ liệu (mục~4.3.3), đảm bảo StockBalance luôn khớp tuyệt đối với tổng Ledger.

\subsection{Script DDL (Định nghĩa bảng)}

Dưới đây là các bảng cốt lõi nhất, trích từ file \texttt{schema.sql} đầy đủ (17 bảng). Quy ước chung: khóa chính dùng \texttt{UUID} (qua \texttt{gen\_random\_uuid()}) thay vì số tự tăng, để phù hợp mô hình phân tán multi-tenant và không lộ số lượng bản ghi qua ID.

\begin{lstlisting}
CREATE EXTENSION IF NOT EXISTS "pgcrypto";

CREATE TABLE tenants (
    id          UUID PRIMARY KEY DEFAULT gen_random_uuid(),
    name        VARCHAR(255) NOT NULL,
    status      VARCHAR(20)  NOT NULL DEFAULT 'active',
    created_at  TIMESTAMPTZ  NOT NULL DEFAULT now()
);

CREATE TABLE branches (
    id          UUID PRIMARY KEY DEFAULT gen_random_uuid(),
    tenant_id   UUID NOT NULL REFERENCES tenants(id),
    name        VARCHAR(255) NOT NULL,
    address     VARCHAR(500),
    is_active   BOOLEAN NOT NULL DEFAULT true,
    created_at  TIMESTAMPTZ NOT NULL DEFAULT now()
);

CREATE TABLE users (
    id             UUID PRIMARY KEY DEFAULT gen_random_uuid(),
    tenant_id      UUID NOT NULL REFERENCES tenants(id),
    email          VARCHAR(255) NOT NULL,
    password_hash  VARCHAR(255) NOT NULL,
    role           VARCHAR(20) NOT NULL
                   CHECK (role IN ('Owner','Staff','Cashier')),
    created_at     TIMESTAMPTZ NOT NULL DEFAULT now(),
    UNIQUE (tenant_id, email)
);
\end{lstlisting}

\begin{lstlisting}
CREATE TABLE product_variants (
    id               UUID PRIMARY KEY DEFAULT gen_random_uuid(),
    tenant_id        UUID NOT NULL REFERENCES tenants(id),
    product_id       UUID NOT NULL REFERENCES products(id),
    sku_code         VARCHAR(64) NOT NULL,
    attributes       JSONB,
    retail_price     NUMERIC(14,2) NOT NULL DEFAULT 0,
    wholesale_price  NUMERIC(14,2),
    created_at       TIMESTAMPTZ NOT NULL DEFAULT now(),
    UNIQUE (tenant_id, sku_code)
);

-- "Read model" tồn kho: dẫn xuất từ Ledger, xem mục 4.2.4
CREATE TABLE stock_balances (
    sku_id      UUID NOT NULL REFERENCES product_variants(id),
    branch_id   UUID NOT NULL REFERENCES branches(id),
    on_hand     NUMERIC(14,3) NOT NULL DEFAULT 0,
    reserved    NUMERIC(14,3) NOT NULL DEFAULT 0,
    avg_cost    NUMERIC(14,4) NOT NULL DEFAULT 0,
    updated_at  TIMESTAMPTZ   NOT NULL DEFAULT now(),
    PRIMARY KEY (sku_id, branch_id),
    CONSTRAINT chk_available_non_negative CHECK (on_hand - reserved >= 0)
);
\end{lstlisting}

\begin{lstlisting}
-- Sổ cái append-only -- nguồn sự thật duy nhất về tồn kho
CREATE TABLE inventory_transactions (
    id             UUID PRIMARY KEY DEFAULT gen_random_uuid(),
    tenant_id      UUID NOT NULL REFERENCES tenants(id),
    branch_id      UUID NOT NULL REFERENCES branches(id),
    sku_id         UUID NOT NULL REFERENCES product_variants(id),
    type           VARCHAR(3) NOT NULL CHECK (type IN ('IN','OUT')),
    quantity       NUMERIC(14,3) NOT NULL CHECK (quantity > 0),
    balance_after  NUMERIC(14,3) NOT NULL,
    reference_id   UUID,
    created_at     TIMESTAMPTZ NOT NULL DEFAULT now()
);

CREATE TABLE orders (
    id          UUID PRIMARY KEY DEFAULT gen_random_uuid(),
    tenant_id   UUID NOT NULL REFERENCES tenants(id),
    branch_id   UUID NOT NULL REFERENCES branches(id),
    channel     VARCHAR(20) NOT NULL,
    status      VARCHAR(20) NOT NULL DEFAULT 'Draft'
                CHECK (status IN
                    ('Draft','Reserved','Confirmed','Completed','Cancelled')),
    created_at  TIMESTAMPTZ NOT NULL DEFAULT now()
);

CREATE TABLE order_items (
    id             UUID PRIMARY KEY DEFAULT gen_random_uuid(),
    order_id       UUID NOT NULL REFERENCES orders(id),
    sku_id         UUID NOT NULL REFERENCES product_variants(id),
    quantity       NUMERIC(14,3) NOT NULL CHECK (quantity > 0),
    selling_price  NUMERIC(14,2) NOT NULL,
    cost_price     NUMERIC(14,4) NOT NULL   -- snapshot bất biến, xem BR-04
);
\end{lstlisting}

\textit{9 bảng còn lại (Category, Brand, Product, Barcode, PurchaseReceipt, PurchaseReceiptItem, StockTransfer, Stocktake, StocktakeItem) theo đúng cấu trúc tương tự, xem đầy đủ trong file \texttt{schema.sql} đính kèm.}

\subsection{Cơ chế Append-only cho Sổ cái Tồn kho}
\label{sec:append-only}
Đây là phần triển khai quan trọng nhất, hiện thực hóa BR-01 và NFR-SEC-03: bảng \texttt{inventory\_transactions} tuyệt đối không cho phép \texttt{UPDATE}/\texttt{DELETE}, kể cả khi thao tác SQL trực tiếp (không qua Backend).

Nhóm triển khai bằng \textbf{TRIGGER} (thay vì RULE), vì trigger cho phép chủ động \texttt{RAISE EXCEPTION} — báo lỗi tường minh và rollback giao dịch, thay vì âm thầm bỏ qua thao tác như RULE:

\begin{lstlisting}
CREATE OR REPLACE FUNCTION fn_block_ledger_modification()
RETURNS TRIGGER AS $$
BEGIN
    RAISE EXCEPTION
      'inventory_transactions is append-only. % is not allowed. '
      'Use a compensating entry instead.', TG_OP;
END;
$$ LANGUAGE plpgsql;

CREATE TRIGGER trg_block_update_inventory_transaction
    BEFORE UPDATE ON inventory_transactions
    FOR EACH ROW EXECUTE FUNCTION fn_block_ledger_modification();

CREATE TRIGGER trg_block_delete_inventory_transaction
    BEFORE DELETE ON inventory_transactions
    FOR EACH ROW EXECUTE FUNCTION fn_block_ledger_modification();
\end{lstlisting}

Khi có ai cố tình chạy \texttt{UPDATE}/\texttt{DELETE} trên bảng này, PostgreSQL trả lỗi rõ ràng và hủy toàn bộ giao dịch (rollback), không có dữ liệu nào bị thay đổi. Mọi điều chỉnh sai sót phải thực hiện bằng cách chèn thêm một dòng bù trừ mới (compensating entry) — đúng tinh thần "chỉ ghi thêm, không sửa" của sổ cái kế toán truyền thống. Bằng chứng kiểm thử cơ chế này được trình bày tại mục~4.3.2.

\subsection{Kiểm soát đồng thời (Concurrency Control)}
\label{sec:concurrency-control}
Đây là cơ chế giải quyết trực tiếp vấn đề cốt lõi của đề tài — chống bán vượt tồn kho (NFR-PERF-02). PostgreSQL mặc định chạy ở mức cô lập giao dịch \texttt{READ COMMITTED}, không đủ để tự chống race-condition khi nhiều giao dịch đồng thời đọc rồi ghi lại cùng một dòng. Nhóm áp dụng \textbf{Pessimistic Locking} bằng \texttt{SELECT...FOR UPDATE} để khóa đúng dòng \texttt{stock\_balances} đang bị tranh chấp trước khi kiểm tra và cập nhật:

\begin{lstlisting}
BEGIN;

-- Khoa dong StockBalance dang bi tranh chap (block cac transaction khac
-- dang cho khoa tren cung dong nay cho den khi COMMIT/ROLLBACK)
SELECT on_hand, reserved
FROM stock_balances
WHERE sku_id = :sku_id AND branch_id = :branch_id
FOR UPDATE;

-- Tang lai o tang ung dung: available = on_hand - reserved
-- Neu available < so luong dat mua -> ROLLBACK, tra loi 409 Conflict

UPDATE stock_balances
SET reserved = reserved + :requested_qty,
    updated_at = now()
WHERE sku_id = :sku_id AND branch_id = :branch_id;

COMMIT;
\end{lstlisting}

Với cơ chế này, nếu 50 giao dịch đồng thời cùng cố gắng \texttt{SELECT...FOR UPDATE} trên một dòng \texttt{stock\_balances} chỉ còn 1 đơn vị tồn, PostgreSQL sẽ tuần tự hóa (serialize) các giao dịch: chỉ giao dịch đầu tiên lấy được khóa và commit thành công, 49 giao dịch còn lại phải chờ, và khi lấy được khóa sẽ thấy \texttt{reserved} đã cập nhật nên tự động bị từ chối ở bước kiểm tra \texttt{available} — không có khả năng 2 giao dịch cùng đọc được giá trị cũ và cùng ghi đè lẫn nhau (lost update).

\subsection{Chỉ mục (Index) tối ưu}
Theo NFR-TENANT-02, các chỉ mục composite sau được tạo để duy trì hiệu năng truy vấn khi bảng ledger đạt quy mô hàng triệu dòng:

\begin{lstlisting}
CREATE INDEX idx_ledger_tenant_created
    ON inventory_transactions (tenant_id, created_at);

CREATE INDEX idx_ledger_tenant_sku
    ON inventory_transactions (tenant_id, sku_id);

CREATE INDEX idx_orders_tenant_created
    ON orders (tenant_id, created_at);
\end{lstlisting}

\textbf{Lý do chọn thứ tự cột:} TenantId luôn đứng đầu vì mọi truy vấn của hệ thống multi-tenant đều lọc theo Tenant trước tiên (áp dụng Global Query Filter) — đặt TenantId đầu composite index giúp PostgreSQL thu hẹp phạm vi quét ngay từ bước đầu (index prefix), trước khi lọc tiếp theo CreatedAt hoặc SKUId. Nếu đảo ngược thứ tự (SKUId trước TenantId), index sẽ kém hiệu quả hơn cho các báo cáo lọc theo khoảng thời gian trên toàn bộ Tenant.

\subsection{Dữ liệu mẫu (Seed Data)}
Để minh họa schema hoạt động đúng, nhóm xây dựng bộ dữ liệu mẫu gồm 2--3 Tenant độc lập, mỗi Tenant có User theo cả 3 vai trò, sản phẩm/SKU với tồn kho ban đầu, và đơn hàng mẫu ở đủ các trạng thái. Ví dụ một đoạn seed data minh họa kịch bản nhập hàng 2 đợt với giá khác nhau — dùng để kiểm chứng công thức giá vốn bình quân ở mục~4.3.3:

\begin{lstlisting}
-- Dot nhap 1: 100 don vi, gia 50.000d -> AvgCost = 50.000
INSERT INTO inventory_transactions
    (tenant_id, branch_id, sku_id, type, quantity, balance_after, created_at)
VALUES
    (:tenant_id, :branch_id, :sku_id, 'IN', 100, 100, '2026-08-01');
UPDATE stock_balances SET on_hand = 100, avg_cost = 50000
    WHERE sku_id = :sku_id AND branch_id = :branch_id;

-- Dot nhap 2: 50 don vi, gia 56.000d
-- AvgCost moi = (100*50000 + 50*56000) / 150 = 52.000
INSERT INTO inventory_transactions
    (tenant_id, branch_id, sku_id, type, quantity, balance_after, created_at)
VALUES
    (:tenant_id, :branch_id, :sku_id, 'IN', 50, 150, '2026-08-10');
UPDATE stock_balances SET on_hand = 150, avg_cost = 52000
    WHERE sku_id = :sku_id AND branch_id = :branch_id;
\end{lstlisting}

\textit{Lưu ý: câu lệnh \texttt{UPDATE stock\_balances} ở trên hợp lệ vì bảng này KHÔNG bị ràng buộc append-only (chỉ \texttt{inventory\_transactions} mới bị chặn UPDATE/DELETE — xem mục~4.2.4).}
